\honest{Occam's Razor}{Eliezer Yudkowsky}{2007}

The more complex an explanation, the more evidence you need just to find it. How do we measure complexity? Occam's razor is often put as ``The simplest explanation that fits the facts,'' and Heinlein, I say, replied that the simplest explanation is ``The lady down the street is a witch; she did it.''

So the length of an English sentence is a poor measure. Words are labels for things the listener already knows, and a short label, ``Fnord!'', does not make a claim simple. ``Witch'' and ``anger'' sound simple only because we do not see the brain's machinery behind them. To a human, Maxwell's equations take longer to explain than Thor. But a program that simulates Maxwell's equations is ``enormously'' shorter than one that simulates a mind, ``as it turns out.''

Solomonoff induction measures complexity by the length of the shortest program that produces a description. If you dislike its Turing machines, I say, ``there's only a constant complexity penalty'' for using another language. \nb{A language can hide complexity in a label, as English does. A universal machine with Thor as a one-bit instruction is still universal, and on it Thor is simpler than Maxwell; the invariance theorem permits this, because its constant can be as large as the hypotheses being compared. So the formalism alone does not rank Maxwell above Thor. In 2015 Leike and Hutter called the choice of machine ``a big open question.''}

In the version I prefer, programs assign probabilities to data. A program fixed to show HTTHHT fits that sequence 64 times better than a fair coin, so fit must be paid for with complexity: each bit stored in the program must buy at least a doubling of fit, or ``you will sooner or later decide that all fair coins are fixed.'' Solomonoff's predictor weights every program by 2 to the minus its length and by its fit, and it is uncomputable. \nb{This is the post's answer to how complexity can be measured. Apart from the comparison of Maxwell and Thor, the post gives no way to estimate it for a real explanation.} Minimum Message Length is nearly the same: send a code, then the data in that code, and prefer the shortest total.

Now the witch. To tell a friend the sequence 0101010101, you would still have to send all ten digits after saying ``witch.'' Witchcraft, like phlogiston, permits everything, so it compresses nothing; it is ``a useless prologue.'' The trick was in the word ``it.''

Why the witch still seems to explain the data, I promise to say later.
